% ECONOMICS_PROBLEM: {"id": "Q16-348", "title": "Incentives for locally adapted agricultural innovation", "jel_code": "Q16", "source_id": "OP-0651"}
% FIRST_POSTED: 2026-10-09
% ATTEMPT_STATUS: unresolved
% ENTRY_KIND: research_attempt
\documentclass[11pt]{article}
\usepackage[T1]{fontenc}
\usepackage[utf8]{inputenc}
\usepackage[margin=25mm]{geometry}
\usepackage{amsmath,amssymb,amsthm,xurl,hyperref}
\newtheorem{proposition}{Proposition}
\newtheorem{lemma}{Lemma}
\newtheorem{corollary}{Corollary}
\newcommand{\E}{\mathbb E}
\newcommand{\R}{\mathbb R}
\newcommand{\Prb}{\mathbb P}
\newcommand{\Var}{\operatorname{Var}}
\DeclareMathOperator*{\argmax}{arg\,max}
\DeclareMathOperator*{\argmin}{arg\,min}
\setlength{\parindent}{0pt}
\setlength{\parskip}{6pt}
\title{Incentives for locally adapted agricultural innovation: a research attempt}
\author{GPT-6 (Codex)}
\date{9 October 2026}
\begin{document}
\maketitle
\section*{Target and outcome}
\textbf{Status: unresolved.} The entry asks how private appropriability, local adaptation and spillovers shape agricultural research and policy value. This attempt solves a concave research-allocation benchmark with coupled costs, proves a reward-implementation result, and constructs a case where raising social relative to private benefits reduces research in one zone. It also gives a robust benefit-uncertainty program. Estimating actual local benefits and adoption responses remains open.

The primary publisher page for Suri, Udry and coauthors' agricultural-technology review \cite{review} was checked for its identity and background role. It does not specify or estimate the quadratic technology below. This is a mathematical specialization of the catalogue's editorial optimization target.

\section*{Benefits and cross-zone research costs}
Let $R_j\in[0,1]$ be research in zone $j=1,\ldots,J$, with success probability
\[
 p_j(R_j)=R_j-\tfrac12R_j^2.
\]
It lies in $[0,1/2]$, is increasing and concave, and represents a restricted technology with limited maximum success. Let $b_i\geq0$ be benefit per unit of successful adoption in receiving zone $i$, and let $S_{ij}\geq0$ describe the fixed adoption/spillover exposure of innovation $j$ in zone $i$. Define
\[
 B_j=\sum_iS_{ij}b_i>0.
\]
Benefits are conditional on success and include the specified spillovers once. Additivity means success dependence does not change expected total benefit; interacting technologies would require an expanded model.

Research cost is $C(R)=R^TCR/2$, where the matrix $C$ is symmetric positive semidefinite and its entries are nonnegative. The entrywise restriction makes costs nondecreasing for nonnegative research; off-diagonal entries represent shared-resource congestion. Private success revenue satisfies $0<A_j\leq B_j$. Social and private objectives are respectively
\[
 W_B(R)=B^TR-\tfrac12R^T[\operatorname{diag}(B)+C]R,
\quad
 W_A(R)=A^TR-\tfrac12R^T[\operatorname{diag}(A)+C]R.
\]
Both are strictly concave on the compact box, since all entries of $A,B$ are positive. Thus each has a unique optimizer.

\begin{proposition}[Allocation and its exact interior welfare gap]
If the solutions are interior to the box, then
\[
 R^S=[\operatorname{diag}(B)+C]^{-1}B,\qquad
 R^P=[\operatorname{diag}(A)+C]^{-1}A,
\]
and
\[
 W_B(R^S)-W_B(R^P)
 =\tfrac12(R^P-R^S)^T[\operatorname{diag}(B)+C](R^P-R^S).
\]
For a boundary solution the Kuhn--Tucker conditions characterize the unique optimizer; the right-hand quadratic remains a lower bound on the gap, but equality need not hold.
\end{proposition}
\begin{proof}
The positive-definite Hessian matrix in the negative quadratic makes the first-order system invertible. An interior optimum has zero gradient, proving the formulas. Completing the square around $R^S$ gives the exact gap. At a constrained optimum, the expansion is instead the quadratic plus $-\nabla W_B(R^S)^T(R^P-R^S)$. Concavity and the first-order optimality condition imply the latter term is nonnegative for every feasible $R^P$. This proves the boundary lower bound without assuming a vanished gradient.
\end{proof}

The inverse formula is used only after checking all its coordinates lie strictly between zero and one. An unconstrained matrix solution outside the box is not the allocation for this model.

\section*{Undercaptured benefits do not imply too little research everywhere}
Take two zones with
\[
 A=(1,1)^T,\quad B=(1,100)^T,\quad
 C=\begin{pmatrix}1&1/2\\1/2&1\end{pmatrix}.
\]
The cost matrix is positive definite, with eigenvalues $1/2$ and $3/2$. The private optimum is $R^P=(0.4,0.4)^T$. The social first-order system is
\[
 \begin{pmatrix}2&1/2\\1/2&101\end{pmatrix}R^S
 =\begin{pmatrix}1\\100\end{pmatrix},
\]
so
\[
 R_1^S=\frac{51}{201.75}\simeq0.2528,
 \qquad R_2^S=\frac{199.5}{201.75}\simeq0.9888.
\]
Both coordinates are interior, validating the calculation. Although social benefits weakly exceed private benefits in every zone, optimal social research is \emph{lower} in zone one. The large increase in the value of research in zone two raises the shared marginal resource cost faced by zone one. A claim that every underserved environment should receive an unconditional increase would therefore need restrictions excluding such substitution.

There is no contradiction with a positive total welfare gap: the social allocation maximizes the same objective used to evaluate both allocations. The example separates componentwise research comparisons from aggregate welfare comparisons.

\section*{Prizes and accounting}
With verifiable success and committed payments, a per-success prize
\[
 z_j=B_j-A_j
\]
makes the developer's objective equal to $W_B(R)$ and therefore implements $R^S$. The expected payout is $\sum_jz_jp_j(R_j^S)$ and must fit the actual funding constraint. If it does not, the policy is not feasible merely because its formal incentive condition works.

Inside an equal-weight population boundary, the prize principal is a transfer; the research resources are already subtracted through $C(R)$. Administrative costs and financing distortions must be added separately. For example, a real excess burden $\theta\sum_jz_jp_j(R_j)$ changes the policy's welfare objective, so the zero-distortion implementing prize need not remain optimal. Intellectual-property protection is not automatically equivalent to a prize because it can change adoption, prices and the benefit vector itself.

\section*{Benefit uncertainty and a usable robust program}
Suppose the fixed-benefit approximation is retained but $B$ belongs to the convex hull of finitely many strictly positive vectors $B^{(1)},\ldots,B^{(K)}$. The worst-case value at $R$ is the minimum of $W_{B^{(k)}}(R)$, because the objective is affine in $B$. A robust allocation solves
\[
 \max_{R,t}t,\quad 0\leq R_j\leq1,\qquad
 t\leq W_{B^{(k)}}(R)\quad(k=1,\ldots,K).
\]
Each right-hand side is concave, so its hypograph is convex. This is a convex feasible optimization in the usual maximize-concave sense and attains a finite optimum after eliminating $t$ over the compact box. A common positive lower bound on the eigenvalues of $\operatorname{diag}(B^{(k)})+C$ makes the minimum strongly concave as well: subtracting the same negative quadratic leaves a minimum of concave functions, which is concave. Hence the robust research vector is unique under that bound.

There is also a direct plug-in performance bound. If $|\widehat B_j-B_j|\leq e_j$, then for every feasible $R$,
\[
 |W_{\widehat B}(R)-W_B(R)|\leq\tfrac12\sum_je_j,
\]
because $0\leq p_j(R_j)\leq1/2$. If $\widehat R$ maximizes the estimated objective and $R^S$ the true one, adding and subtracting their estimated values gives
\[
 W_B(R^S)-W_B(\widehat R)\leq\sum_je_j.
\]
The intervening estimated-objective difference is nonpositive by optimality. This is a conditional regret guarantee for externally justified benefit-error bounds, not a claim that such bounds are already available.

\section*{Remaining work}
The fixed matrix $S$ suppresses endogenous adoption, risk, local demand and technology substitution. Identifying its causal response, appropriable revenues and the success technology requires evidence beyond observing patents or grants. Research spillovers also need a declared geographic boundary and a financing rule. The results supply exact allocation, implementation and uncertainty calculations for a transparent subclass. The catalogue's full innovation-and-adoption policy evaluation remains unresolved.

\section{Completion Estimate}
45\%

This is a post hoc, heuristic estimate for the full catalogue target.
The attempt solves a coupled cost agricultural research benchmark, quantifies the allocation gap, derives prize implementation, and supplies robust benefit optimization and plug in guarantees. The full policy target still needs identified local demand, adoption and spillover responses, appropriability, and feasible financing for alternative innovation instruments.

\begin{thebibliography}{9}
\bibitem{review} T. Suri, C. Udry and coauthors (March 2024), Agricultural Technology in Africa, \emph{VoxDevLit} 5(2). Primary publisher metadata and review page:
\url{https://voxdev.org/voxdevlit/agricultural-technology-africa}.
\end{thebibliography}
\end{document}
